% Chapter 2 — Overall Description
\section{Overall Description}

\subsection{Product perspective}
% TODO: Scenarios and further details on the shared phenomena. Domain model:
%       UML class diagrams and state diagrams for the key entities
%       (Session, Car, Driver, Lap, Sector, TelemetrySample, Position,
%       WeatherReading, PitStop, RaceControlEvent, Consumer, Subscription, ...).

\subsection{Product functions}
% TODO: The most important (high-level) requirements: ingest OpenF1 data, store
%       persistently, let consumers subscribe to a subset of information, expose
%       it via GUI and API, replay historical sessions.

\subsection{User characteristics}
% TODO: Who the users are and what clarifies their needs: racing enthusiasts
%       (replay, gaps, events) vs. third-party developers (telemetry, analyses),
%       differing levels of technical expertise.

\subsection{Assumptions, dependencies and constraints}
% TODO: Domain assumptions [A-1..] with one-line rationale. E.g. OpenF1 is the
%       sole external source; historical data free / live out of scope; ~3.7
%       samples/sec/car; OpenF1 data semantics. Dependencies: OpenF1 availability
%       and rate limits, network, API terms.
